\section{Exact formulation}
\label{sec:exact}

\begin{inbrief}
To measure how far the heuristic is from the optimum of (O1), the
allocation of the sub-chains is also solved exactly. A constraint program
chooses the cores and the LET periods; schedulability is added lazily, as
cuts derived from unschedulable cores. Every optimal allocation it returns
is checked again by the implementation of the heuristic.
\end{inbrief}

\paragraph{Scope.} The exact model takes the sub-chains of step 1 as
given, since partitioning does not depend on the allocation rule, and
solves the rest of (O1): the allocation, the LET periods and the
schedulability of every core, under (C1)--(C7). It uses the same WCETs,
utilizations (in parts per million of a core, rounded up) and LET tasks
as the heuristic, so the two are compared on exactly the same problem.
Contention is left out, as it would make the WCETs depend on the number of
cores per ECU. The candidate cores of an ECU are its initial cores, then
cores of its added type up to two more than the larger of its initial
count and the largest number any heuristic uses there (within
$M^{\max}_k$). The optimum is the optimum over these candidates. A
solution that uses every candidate of an ECU is reported as possibly
restricted; this happened once (\cref{sec:rq1}). To check that two spare
cores suffice, all 240 instances of \cref{sec:rq1} were solved again with
four: no allocation with fewer cores was found on any of them, none was
restricted, and every optimum proven in both runs was the same.

\subsection{Variables and constraints}

For every sub-chain $f$ and allowed candidate core $p$ (by (C1)), a Boolean
$x_{f,p}$ says that $f$ runs on $p$, and $u_p$ that $p$ is used. Let
$\mathcal{G}$ be the closure under GCD of the GCDs $g_f$ of the periods of
the sub-chains; for every core, a Boolean $y_{p,g}$, $g\in\mathcal{G}$,
chooses the period of its LET task. With $U_{f,p}$ the utilization of $f$
on $p$ and $C^{\LET}_p$ the WCET of the LET task:
\begin{align}
&\textstyle\sum_{p}x_{f,p}=1, \qquad x_{f,p}\le u_p &&\forall f,p\label{x:assign}\\
&\textstyle\sum_{g\in\mathcal{G}}y_{p,g}=u_p,\qquad
	x_{f,p}\le\sum_{g\mid g_f}y_{p,g} &&\forall f,p\label{x:let}\\
&\textstyle\sum_{f}U_{f,p}\,x_{f,p}+\sum_{g}\ceil{\frac{C^{\LET}_p}{g}}y_{p,g}
	\le\varphi\,u_p &&\forall p\label{x:load}\\
&\textstyle\sum_{f,\,p\in k}\mu_f\,x_{f,p}\le M_k &&\forall k\label{x:memory}\\
&\textstyle\sum_{\xi\in\varrho}z_{\xi,k}\le1,\quad x_{f,p}\le z_{\xi(f),\kappa(p)}
	&&\forall\varrho,k\label{x:replica}
\end{align}
\Cref{x:let} lets a sub-chain run on a core only if the LET period of the
core divides the periods of its tasks. The chosen $g$ then divides every
period on the core, so it is at most their GCD, and the load it reserves is
at least the true LET load: every solution is admissible. Conversely, the
GCD of the periods on a core is in $\mathcal{G}$, so an allocation that is
admissible with the true LET periods stays feasible with $g$ equal to that
GCD, and no allocation is lost. \Cref{x:load} is the cap (C3), \cref{x:memory} the
memory (C4), and \cref{x:replica} keeps the task-chains $\xi$ of a replica
group $\varrho$ on distinct ECUs (C5), with an auxiliary $z_{\xi,k}$ that is
true if $\xi$ has a sub-chain on ECU $k$; the core policy replaces ECUs by
cores. Co-location (C2) is already ensured by partitioning. The objective
is $\min\sum_p u_p$.

\paragraph{Symmetry breaking.} Cores of one ECU and one kind (type, speed,
LET WCET) are interchangeable. Within such a group, $u_{p_{j+1}}\le u_{p_j}$:
cores are used in order. Moreover, with the sub-chains numbered by
decreasing utilization, sub-chain $i$ may only use the first $i+1$ cores of
a group. Any solution can be relabelled to satisfy this, by sorting the
used cores of each group by the smallest index of their sub-chains: the
$j$-th core then holds no sub-chain of index below $j$.

\subsection{Schedulability as lazy constraints}

Constraints (C6) and (C7) are not linear. They are added lazily
(\cref{alg:exact}), in the spirit of logic-based Benders decomposition
\cite{cambazard2004decomposition}. After each solve, every used core is
checked with the response-time analysis and the priority assignment of
\cref{alg:audsley}, which is optimal (\cref{thm:audsley}). For an
unschedulable core with sub-chains $S$, the check removes sub-chains from
$S$ one at a time, as long as the rest stays unschedulable. The result $S'$
is unschedulable, but it becomes schedulable if any of its sub-chains is
removed. The cut
\begin{equation}
\textstyle\sum_{f\in S'}x_{f,q}\le|S'|-1
\label{eq:cut}
\end{equation}
is then added for every candidate core $q$ of the same kind as the failing
one, and the model is solved again, starting from the previous solution.

\begin{algorithm2e}[t]
\caption{\textsf{Exact}$(\Frags)$}
\label{alg:exact}
$\mathcal{K}\gets\varnothing$ \tcp*{cuts}
\Repeat{every used core is schedulable}{
	solve \cref{x:assign,x:let,x:load,x:memory,x:replica} with the cuts $\mathcal{K}$, minimizing $\sum u_p$\;
	\lIf{infeasible}{\Return \emph{infeasible}}
	\ForEach{used core $p$ with sub-chains $S$, unschedulable}{
		$S'\gets S$\;
		\ForEach{$f\in S'$}{\lIf{$S'\setminus\{f\}$ unschedulable}{$S'\gets S'\setminus\{f\}$}}
		add \cref{eq:cut} for $S'$ on every core of the kind of $p$ to $\mathcal{K}$\;
	}
}
\Return the allocation and $\sum u_p$\;
\end{algorithm2e}

\begin{proposition}[Validity of the cuts]
\label{prop:cuts}
No cut removes an allocation in which every core is schedulable.
\Cref{alg:exact} therefore returns an optimal allocation, and the
objective bound of any of its solves is a lower bound of the optimum.
\end{proposition}
\begin{proof}
Schedulability under the optimal assignment is monotone with respect to
removing sub-chains: removing tasks shrinks every higher-priority set and
can only make the LET period longer, so the order that was schedulable
stays schedulable for the remaining tasks. Hence every superset of the
unschedulable set $S'$ is unschedulable on a core of that kind, and the WCETs,
utilizations and LET task depend only on the kind. An allocation that puts
all of $S'$ on such a core is not admissible, which is what \cref{eq:cut}
forbids. The cuts only remove inadmissible allocations, so every solve is
a relaxation of the problem and its bound is a lower bound. The loop ends:
a new cut is never satisfied by the solution that produced it, so no cut
is added twice, and there are finitely many. When the last solve is proved
optimal, its solution is admissible and optimal for a relaxation of the
problem, hence optimal for the problem. When the time limit stops a solve,
the best admissible solution and the bound are reported instead.
\end{proof}

\paragraph{Implementation and checking.} The model is solved with the
CP-SAT solver of OR-Tools \cite{ortools}, with a time limit; if the limit is
reached, the best admissible solution and the bound are reported. The
schedulability checker is a separate implementation of \cref{eq:rta} and
\cref{alg:audsley}. It is first compared with the heuristic on the
heuristic's own mappings. Every exact allocation is then placed again by
the implementation of the heuristic, with all its admissibility checks and
its priority assignment; the experiments stop if it is rejected. None was.
